\song{Dokud se zpívá (Jaromír Nohavica)}
%Mara_2016_09_05

\ve{}
\ch{C}Z Těšína \ch{Emi}vyjíždí \ch{Dmi7}vlaky co čtv\ch{F}rthodi\ch{C}nu, \ch{Emi} \quad \ch{Dmi} \quad \ch{G}\quad\\
\ch{C}včera jsem \ch{Emi}nespal a \ch{Dmi7}ani dnes \ch{F}nespoči\ch{C}nu,\ch{Emi} \quad \ch{Dmi} \quad \ch{G}\quad\\
\ch{F}svatý \ch{G}Medard, můj \ch{C}patron, ťuká si na \ch{Ami}{č}elo, \ch{G}\quad\\
ale \ch{F}dokud se \ch{G}zpívá, \ch{F}ještě se \ch{G}neumře\ch{C}lo, hó\ch{Emi}hó \quad \ch{Dmi} \quad \ch{G}\quad

\ve{}
Ve stánku koupím si housku a slané tyčky,\\
srdce mám pro lásku a hlavu pro písničky,\\
ze školy dobře vím, co by se dělat mělo,\\
ale dokud se zpívá, ještě se neumřelo, hóhó. 

\ve{}
Do alba jízdenek lepím si další jednu,\\
vyjel jsem před chvílí, konec je v nedohlednu,\\
za oknem míhá se život jak leporelo,\\
ale dokud se zpívá, ještě se neumřelo, hóhó. 

\ve{}
Stokrát jsem prohloupil a stokrát platil draze,\\
houpe to, houpe to na housenkové dráze,\\
i kdyby supi se slítali na mé tělo,\\
tak dokud se zpívá, ještě se neumřelo, hóhó. 

\ve{}
Z Těšína vyjíždí vlaky až na kraj světa,\\
zvedl jsem telefon a ptám se: \uv{Lidi, jste tam?}\\
A z veliké dálky do uší mi zaznělo,\\
/: že dokud se zpívá, ještě se neumřelo. :/ 

\esong{}